\documentclass[a4paper,12pt]{article}

\usepackage[a4paper,margin=16mm]{geometry}
\usepackage[german,main=french]{babel}
\usepackage{fontspec}
\usepackage{microtype}

\setmainfont{EBGaramond}
[
  Extension      = .otf,
  UprightFont    = *-Regular,
  ItalicFont     = *-Italic,
  BoldFont       = *-Bold,
  BoldItalicFont = *-BoldItalic,
  Numbers        = Lowercase
]

\usepackage{paracol}
\newcommand\chunks[2]{%
\begin{leftcolumn*}\begin{otherlanguage}{german}%
{#1}%
\end{otherlanguage}\end{leftcolumn*}%
\begin{rightcolumn}\begin{otherlanguage}{french}%
{#2}%
\end{otherlanguage}\end{rightcolumn}%
}

\newcommand\strophe[1]{\subsection*{\normalsize\normalfont\itshape #1}}

\begin{document}\thispagestyle{empty}
\setlength{\columnsep}{8mm}


\begin{paracol}{2}

  \chunks{\section*{Ein Mädchen oder Weibchen}}
         {\section*{Une fille ou une petite femme}}

  \chunks{\noindent\raggedright\small Wolfgang Amadeus Mozart, \textit{Die Zauberflöte}}
         {\noindent\raggedright\small Wolfgang Amadeus Mozart, \textit{La Flûte enchantée}}

  \chunks{\noindent\raggedright\small Libretto von Emanuel Schikaneder (1791)}
         {\noindent\raggedright\small Livret d’Emanuel Schikaneder (1791)}

  \chunks{}
         {\noindent\raggedright\small Texte français : Jean Spenlehauer}

  \chunks{\strophe{Zweiter Aufzug, dreiundzwanzigster Auftritt, Nr.~20, Arie}}
         {\strophe{Deuxième acte, vingt-troisième scène, n\textsuperscript{o}~20, air}}

  \chunks{\noindent\small(Er schlägt das Glockenspiel.)}
         {\noindent\small(Il joue sur son glockenspiel.)}

  \chunks{\begin{verse}
          \textsc{Papageno}\\[\medskipamount]
          Ein Mädchen oder Weibchen wünscht Papageno sich!\\
          O so ein sanftes Täubchen wär’ Seligkeit für mich!\\
          Dann schmeckte mir Trinken und Essen\\
          Dann könnt’ ich mit Fürsten mich messen,\\
          Des Lebens als Weiser mich freun,\\
          Und wie im Elysium sein.\\
          Dann könnt’ ich, mit Fürsten, mich messen,\\
          Des Lebens als Weiser mich freun,\\
          Und wie im Elysium sein, im Elysium sein,\\
          Im Elysium sein!
          \end{verse}}
         {\begin{verse}
          \textsc{Papageno}\\[\medskipamount]
          Une fille ou une petite femme à lui,\\
          c’est ce que désire Papageno !\\
          O une si douce petite colombe serait la félicité pour moi !\\
          Alors j’apprécierais le boire et le manger,\\
          Alors le roi ne serait pas mon cousin,\\
          Je jouirais de la vie en sage,\\
          Je serais comme en paradis.\\
          Alors le roi ne serait pas mon cousin,\\
          Je jouirais de la vie en sage,\\
          Je serais comme en paradis,\\
          En paradis, en paradis.
          \end{verse}}

  \chunks{\begin{verse}
          Ein Mädchen oder Weibchen wünscht Papageno sich!\\
          O so ein sanftes Täubchen wär’ Seligkeit für mich!\\
          Ach kann ich denn keiner von allen\\
          Den reizenden Mädchen gefallen?\\
          Helf’ eine mir nur aus der Not,\\
          Sonst gräm’ ich mich wahrlich zu Tod.\\
          Ach kann ich, denn keiner, gefallen\\
          Helf’ eine mir nur aus der Not,\\
          Sonst gräm’ ich mich wahrlich zu Tod.\\
          Mich wahrlich zu Tod, mich wahrlich zu Tod.
          \end{verse}}
         {\begin{verse}
          Une fille ou une petite femme à lui,\\
          c’est ce que désire Papageno !\\
          O une si douce petite colombe serait la félicité pour moi !\\
          Ah, ne puis-je plaire à aucune\\
          De ces charmantes jeunes filles ?\\
          Il faut qu’une me sorte de ma peine,\\
          Sinon, vraiment, je mourrai de chagrin.\\
          Ah, ne puis-je plaire à aucune\\
          Il faut qu’une me sorte de ma peine,\\
          Sinon, vraiment, je mourrai de chagrin.\\
          Je mourrai de chagrin, je mourrai de chagrin.
          \end{verse}}

  \chunks{\begin{verse}
          Ein Mädchen oder Weibchen wünscht Papageno sich!\\
          O so ein sanftes Täubchen wär’ Seligkeit für mich!\\
          Wird keine mir Liebe gewähren,\\
          So muß mich die Flamme verzehren,\\
          Doch küßt mich ein weiblicher Mund\\
          So bin ich schon wieder gesund,\\
          Doch küßt mich ein weiblicher Mund\\
          Doch küßt mich ein weiblicher Mund\\
          So bin ich schon wieder gesund,\\
          Schon wieder gesund, schon wieder gesund.
          \end{verse}}
         {\begin{verse}
          Une fille ou une petite femme à lui,\\
          c’est ce que désire Papageno !\\
          O une si douce petite colombe serait la félicité pour moi !\\
          Si aucune ne m’accorde son amour,\\
          Alors les flammes me dévoreront,\\
          Mais que m’embrasse une bouche féminine,\\
          Alors je retrouverai la santé.\\
          Mais que m’embrasse une bouche féminine\\
          Mais que m’embrasse une bouche féminine,\\
          Alors je retrouverai la santé,\\
          Retrouverai la santé, retrouverai la santé.
          \end{verse}}
\end{paracol}
\end{document}
